% Editable pure source geometry; analysis figures are numbered by the overlay.
% Deterministic editable paths; see make_electrostatic_tikz.py.
\newcommand{\ThUnequalAnalysisDrawing}{\begin{tikzpicture}[>=Latex,font=\small,x=.9cm,y=.9cm]
\begin{scope}\clip(-3.5,-2.8)rectangle(5.3,2.8);
\draw[DeepBlue,line width=.65pt] plot coordinates{(-1.09326,-0.04475) (-2.24018,-0.61600) (-3.50000,-1.27733)};
\draw[->,DeepBlue,line width=.65pt] (-2.33648,-0.66539)--(-2.51420,-0.75714);
\draw[DeepBlue,line width=.65pt] plot coordinates{(-3.50000,1.27733) (-2.24018,0.61600) (-1.09326,0.04475)};
\draw[->,DeepBlue,line width=.65pt] (-2.09645,0.54254)--(-2.27455,0.63354);
\draw[DeepBlue,line width=.65pt] plot coordinates{(-1.07143,-0.07261) (-1.63606,-0.66267) (-2.18444,-1.26933) (-2.77817,-1.96000) (-3.47244,-2.80000)};
\draw[->,DeepBlue,line width=.65pt] (-2.35731,-1.46698)--(-2.48793,-1.61842);
\draw[DeepBlue,line width=.65pt] plot coordinates{(-3.47244,2.80000) (-2.77817,1.96000) (-2.18444,1.26933) (-1.63606,0.66267) (-1.07143,0.07261)};
\draw[->,DeepBlue,line width=.65pt] (-2.11838,1.19426)--(-2.25049,1.34441);
\draw[DeepBlue,line width=.65pt] plot coordinates{(-1.05164,-0.08939) (-1.40347,-0.71867) (-1.71474,-1.32533) (-2.03966,-2.01600) (-2.38189,-2.80000)};
\draw[->,DeepBlue,line width=.65pt] (-1.77783,-1.45441)--(-1.86392,-1.63493);
\draw[DeepBlue,line width=.65pt] plot coordinates{(-2.38189,2.80000) (-2.03966,2.01600) (-1.71474,1.32533) (-1.40347,0.71867) (-1.05164,0.08939)};
\draw[->,DeepBlue,line width=.65pt] (-1.64633,1.18690)--(-1.73502,1.36616);
\draw[DeepBlue,line width=.65pt] plot coordinates{(-1.03063,-0.09789) (-1.22883,-0.76533) (-1.38157,-1.38133) (-1.51425,-2.04400) (-1.63555,-2.80000)};
\draw[->,DeepBlue,line width=.65pt] (-1.39329,-1.43285)--(-1.43485,-1.62848);
\draw[DeepBlue,line width=.65pt] plot coordinates{(-1.63555,2.80000) (-1.51425,2.04400) (-1.38157,1.38133) (-1.22883,0.76533) (-1.03063,0.09789)};
\draw[->,DeepBlue,line width=.65pt] (-1.33457,1.17658)--(-1.37973,1.37142);
\draw[DeepBlue,line width=.65pt] plot coordinates{(-1.01032,-0.10267) (-1.06689,-0.80267) (-1.08035,-1.41867) (-1.05389,-2.07200) (-0.98481,-2.80000)};
\draw[->,DeepBlue,line width=.65pt] (-1.07993,-1.48670)--(-1.07469,-1.68663);
\draw[DeepBlue,line width=.65pt] plot coordinates{(-0.98481,2.80000) (-1.05389,2.07200) (-1.08035,1.41867) (-1.06689,0.80267) (-1.01032,0.10267)};
\draw[->,DeepBlue,line width=.65pt] (-1.08082,1.22533)--(-1.08077,1.42533);
\draw[DeepBlue,line width=.65pt] plot coordinates{(-0.98961,-0.10267) (-0.94408,-0.51333) (-0.89573,-0.84933) (-0.83800,-1.15840) (-0.76614,-1.46533) (-0.67613,-1.78267) (-0.57076,-2.10000) (-0.44200,-2.44102) (-0.29155,-2.80000)};
\draw[->,DeepBlue,line width=.65pt] (-0.73844,-1.57165)--(-0.68267,-1.76372);
\draw[DeepBlue,line width=.65pt] plot coordinates{(-0.29155,2.80000) (-0.44200,2.44102) (-0.57076,2.10000) (-0.67613,1.78267) (-0.76614,1.46533) (-0.83800,1.15840) (-0.89573,0.84933) (-0.94408,0.51333) (-0.98961,0.10267)};
\draw[->,DeepBlue,line width=.65pt] (-0.80363,1.31638)--(-0.75504,1.51038);
\draw[DeepBlue,line width=.65pt] plot coordinates{(-0.96821,-0.10114) (-0.81600,-0.56197) (-0.67300,-0.93059) (-0.51900,-1.25698) (-0.33200,-1.58206) (-0.12300,-1.88409) (0.11900,-2.18163) (0.40500,-2.48725) (0.73291,-2.80000)};
\draw[->,DeepBlue,line width=.65pt] (-0.28212,-1.66108)--(-0.16782,-1.82520);
\draw[DeepBlue,line width=.65pt] plot coordinates{(0.73291,2.80000) (0.40500,2.48725) (0.11900,2.18163) (-0.12300,1.88409) (-0.33200,1.58206) (-0.51900,1.25698) (-0.67300,0.93059) (-0.81600,0.56197) (-0.96821,0.10114)};
\draw[->,DeepBlue,line width=.65pt] (-0.44634,1.39318)--(-0.34406,1.56505);
\draw[DeepBlue,line width=.65pt] plot coordinates{(-0.95274,-0.09333) (-0.75000,-0.48386) (-0.58500,-0.76872) (-0.40900,-1.02886) (-0.23300,-1.24617) (-0.03500,-1.44783) (0.17400,-1.62125) (0.40500,-1.77541) (0.64700,-1.90306) (0.98800,-2.03655) (1.36200,-2.13546) (1.73600,-2.19995) (2.75900,-2.31885) (3.33100,-2.41035) (4.04600,-2.57019) (4.86710,-2.80000)};
\draw[->,DeepBlue,line width=.65pt] (1.45012,-2.15413)--(1.64726,-2.18780);
\draw[DeepBlue,line width=.65pt] plot coordinates{(4.86710,2.80000) (4.04600,2.57019) (3.33100,2.41035) (2.75900,2.31885) (1.73600,2.19995) (1.36200,2.13546) (0.98800,2.03655) (0.64700,1.90306) (0.40500,1.77541) (0.17400,1.62125) (-0.03500,1.44783) (-0.23300,1.24617) (-0.40900,1.02886) (-0.58500,0.76872) (-0.75000,0.48386) (-0.95274,0.09333)};
\draw[->,DeepBlue,line width=.65pt] (0.76992,1.95850)--(0.95740,2.02816);
\draw[DeepBlue,line width=.65pt] plot coordinates{(-0.94062,-0.08707) (-0.68400,-0.45420) (-0.47500,-0.72226) (-0.27700,-0.93692) (-0.07900,-1.11072) (0.14100,-1.25881) (0.37200,-1.36782) (0.60300,-1.43220) (0.82300,-1.45182) (0.95500,-1.44232) (1.08700,-1.41414) (1.19935,-1.37200) (1.30700,-1.31035) (1.39500,-1.23638) (1.46646,-1.14800) (1.51500,-1.05467) (1.54431,-0.95200) (1.55273,-0.84933) (1.54275,-0.74667) (1.51243,-0.63467) (1.46272,-0.52267) (1.40060,-0.42000) (1.31511,-0.30800) (1.07600,-0.06750)};
\draw[->,DeepBlue,line width=.65pt] (0.75608,-1.45474)--(0.95592,-1.44687);
\draw[DeepBlue,line width=.65pt] plot coordinates{(1.07600,0.06750) (1.31511,0.30800) (1.40060,0.42000) (1.46272,0.52267) (1.51243,0.63467) (1.54275,0.74667) (1.55273,0.84933) (1.54431,0.95200) (1.51500,1.05467) (1.46646,1.14800) (1.39500,1.23638) (1.30700,1.31035) (1.19935,1.37200) (1.08700,1.41414) (0.95500,1.44232) (0.82300,1.45182) (0.60300,1.43220) (0.37200,1.36782) (0.14100,1.25881) (-0.07900,1.11072) (-0.27700,0.93692) (-0.47500,0.72226) (-0.68400,0.45420) (-0.94062,0.08707)};
\draw[->,DeepBlue,line width=.65pt] (0.26227,1.32525)--(0.44873,1.39758);
\draw[DeepBlue,line width=.65pt] plot coordinates{(-0.92904,-0.07725) (-0.48600,-0.53920) (-0.31000,-0.69569) (-0.14500,-0.81996) (0.02000,-0.91965) (0.18500,-0.99318) (0.35000,-1.03917) (0.50400,-1.05511) (0.66900,-1.03841) (0.81200,-0.98681) (0.93300,-0.90057) (1.02337,-0.78400) (1.07836,-0.64400) (1.09755,-0.48533) (1.08192,-0.30800) (1.03170,-0.10267)};
\draw[->,DeepBlue,line width=.65pt] (0.23562,-1.01525)--(0.43138,-1.05619);
\draw[DeepBlue,line width=.65pt] plot coordinates{(1.03170,0.10267) (1.08192,0.30800) (1.09755,0.48533) (1.07836,0.64400) (1.02337,0.78400) (0.93300,0.90057) (0.81200,0.98681) (0.66900,1.03841) (0.50400,1.05511) (0.35000,1.03917) (0.18500,0.99318) (0.02000,0.91965) (-0.14500,0.81996) (-0.31000,0.69569) (-0.48600,0.53920) (-0.92904,0.07725)};
\draw[->,DeepBlue,line width=.65pt] (-0.04850,0.88502)--(0.12988,0.97547);
\draw[DeepBlue,line width=.65pt] plot coordinates{(-0.91564,-0.06588) (-0.47500,-0.39899) (-0.17800,-0.58711) (-0.03500,-0.65735) (0.09700,-0.70678) (0.22900,-0.73896) (0.35000,-0.75098) (0.47100,-0.74326) (0.59200,-0.71102) (0.69100,-0.66020) (0.77900,-0.58778) (0.85268,-0.49467) (0.90649,-0.39200) (0.95066,-0.26133) (0.98379,-0.10267)};
\draw[->,DeepBlue,line width=.65pt] (0.09853,-0.71208)--(0.29432,-0.75292);
\draw[DeepBlue,line width=.65pt] plot coordinates{(0.98379,0.10267) (0.95066,0.26133) (0.90649,0.39200) (0.85268,0.49467) (0.77900,0.58778) (0.69100,0.66020) (0.59200,0.71102) (0.47100,0.74326) (0.35000,0.75098) (0.22900,0.73896) (0.09700,0.70678) (-0.03500,0.65735) (-0.17800,0.58711) (-0.47500,0.39899) (-0.91564,0.06588)};
\draw[->,DeepBlue,line width=.65pt] (-0.13209,0.61469)--(0.05109,0.69499);
\draw[DeepBlue,line width=.65pt] plot coordinates{(-0.90400,-0.04646) (-0.46400,-0.25397) (-0.17800,-0.37082) (0.05300,-0.43854) (0.26200,-0.46458) (0.44900,-0.44619) (0.53700,-0.41944) (0.62500,-0.37818) (0.77900,-0.26319) (0.93927,-0.07998)};
\draw[->,DeepBlue,line width=.65pt] (0.04182,-0.44003)--(0.23980,-0.46839);
\draw[DeepBlue,line width=.65pt] plot coordinates{(0.93927,0.07998) (0.77900,0.26319) (0.62500,0.37818) (0.53700,0.41944) (0.44900,0.44619) (0.26200,0.46458) (0.05300,0.43854) (-0.17800,0.37082) (-0.46400,0.25397) (-0.90400,0.04646)};
\draw[->,DeepBlue,line width=.65pt] (-0.16710,0.37712)--(0.02438,0.43488);
\draw[->,TextGray](-3.5,0)--(5.3,0);
\end{scope}
\fill[DeepBlue](-1,0)circle(3pt)node[below=4pt]{$+Q$};\fill(1,0)circle(3pt)node[below=4pt]{$-q$};
\fill[Gold](3,0)circle(2pt);\draw[Gold](3,.1)--(4.2,1.35)node[above,align=center]{$E=0$\\$x_0=3a_0$};
\draw[->,Gold](3.12,0)--(3.55,0);\draw[->,Gold](2.88,0)--(2.45,0);
\draw[->,Gold](3,.85)--(3,.4);\draw[->,Gold](3,-.85)--(3,-.4);
\node[below]at(.5,-3.05){图示 $Q=4q,\ a_0=1$；蓝线方向由 $+Q$ 出发，金色箭头示零场点的局部鞍形结构};
\end{tikzpicture}}
\newcommand{\ThDipoleAnalysisDrawing}{\begin{tikzpicture}[>=Latex,font=\small,x=.85cm,y=.85cm]
\draw[->](-.9,0)--(3.3,0)node[right]{$\vect p$};\fill(-.5,0)circle(2pt)node[below]{$-q$};\fill[DeepBlue](.5,0)circle(2pt)node[below]{$+q$};
\coordinate(P)at({2.5*cos(40)},{2.5*sin(40)});\draw[dashed](0,0)--(P)node[midway,above left]{$r$};\fill(P)circle(1.5pt)node[below right]{$P$};
\draw(.8,0)arc(0:40:.8)node[right]{$\theta$};
\draw[->,DeepBlue,thick](P)--++({cos(40)},{sin(40)})node[right]{$E_r$};
\draw[->,Gold,thick](P)--++({-.41955*sin(40)},{.41955*cos(40)})node[above left]{$E_\theta$};
\draw[->,thick](P)--++({cos(40)-.41955*sin(40)},{sin(40)+.41955*cos(40)})node[above]{$\vect E$};
\node[below]at(1.2,-.8){远场分解：$E_r/E_\theta=2\cot\theta$};
\begin{scope}[shift={(7,0)}]
\foreach\C in {1,2,3}{
 \draw[DeepBlue,thick,domain=1:179,samples=140]plot({\C*sin(\x)*sin(\x)*cos(\x)},{\C*sin(\x)*sin(\x)*sin(\x)});
 \draw[DeepBlue,thick,domain=359:181,samples=140]plot({\C*sin(\x)*sin(\x)*cos(\x)},{\C*sin(\x)*sin(\x)*sin(\x)});
 \draw[->,DeepBlue] (.12,\C)--(-.12,\C);\draw[->,DeepBlue](.12,-\C)--(-.12,-\C);
}
\fill(0,0)circle(2pt);\draw[->,thick](0,0)--(.7,0)node[right]{$\vect p$};
\node[above]at(0,3.4){点偶极子：$r=C\sin^2\theta$};
\end{scope}\end{tikzpicture}}
\newcommand{\ThTwoLinesAnalysisDrawing}{\begin{tikzpicture}[>=Latex,font=\small,x=1cm,y=1cm]
\begin{scope}\clip(-4,-2.7)rectangle(4,2.7);
\foreach\C in {1.4,2,3}{\pgfmathsetmacro{\c}{1.5*(\C*\C+1)/(\C*\C-1)}\pgfmathsetmacro{\R}{3*\C/(\C*\C-1)}\draw[DeepBlue](-\c,0)circle(\R);\draw[DeepBlue](\c,0)circle(\R);}
\draw[DeepBlue](0,-2.7)--(0,2.7);
\foreach\d in {-2,-.7,0,.7,2}{\pgfmathsetmacro{\R}{sqrt(2.25+\d*\d)}\draw[TextGray,dashed](0,\d)circle(\R);\draw[->,TextGray](.12,{\d+\R})--(-.12,{\d+\R});\draw[->,TextGray](.12,{\d-\R})--(-.12,{\d-\R});}
\draw[TextGray,dashed](-4,0)--(4,0);\draw[->,TextGray](.2,0)--(-.2,0);
\end{scope}
\fill(-1.5,0)circle(3pt)node[below=4pt]{$-\lambda$};\fill[DeepBlue](1.5,0)circle(3pt)node[below=4pt]{$+\lambda$};
\fill(0,1.5)circle(1.5pt)node[above right]{$P$};\draw(0,1.25)--(-.25,1.25)--(-.25,1.5);
\node[below,align=center]at(0,-3){实线：$r_-/r_+=C$；虚线：$x^2+(y-d)^2=a^2+d^2$\\电荷位置 $x=\pm a$；在 $P=(0,a)$ 可直接检验两族切向正交};
\end{tikzpicture}}
\newcommand{\ThFiniteRodAnalysisDrawing}{\begin{tikzpicture}[>=Latex,font=\small,x=1cm,y=1cm]
\begin{scope}\clip(-4,-2.7)rectangle(4,2.7);
\foreach\A in {1.75,2.2,2.9,3.7}{\pgfmathsetmacro{\B}{sqrt(\A*\A-2.25)}\draw[DeepBlue](0,0)ellipse({\A} and {\B});}
\foreach\eta in {-.92,-.7,-.4,0,.4,.7,.92}{\foreach\sign in {-1,1}{
 \draw[TextGray,dashed,domain=1:3.5,samples=100]plot({1.5*\x*\eta},{\sign*1.5*sqrt((\x*\x-1)*(1-\eta*\eta))});
 \draw[->,TextGray]({1.5*1.3*\eta},{\sign*1.5*sqrt((1.3*1.3-1)*(1-\eta*\eta))})--({1.5*1.4*\eta},{\sign*1.5*sqrt((1.4*1.4-1)*(1-\eta*\eta))});
}}
\end{scope}\draw[very thick](-1.5,0)--(1.5,0);\fill(-1.5,0)circle(2pt)node[below=4pt]{$-a$};\fill(1.5,0)circle(2pt)node[below=4pt]{$a$};
\node[below,align=center]at(0,-3){实线：$r_++r_-=2a\xi$（椭圆）；虚线：$r_--r_+=2a\eta$（双曲线）\\$\lambda>0$ 时箭头沿增大的 $\xi$ 向外，两个端点是共同焦点};
\end{tikzpicture}}
\newcommand{\ThRutherfordAnalysisDrawing}{\begin{tikzpicture}[>=Latex,font=\small,x=.65cm,y=.65cm]
\begin{scope}\clip(-7,-1)rectangle(6.3,6);
\draw[thick,DeepBlue] plot coordinates{(-6.99123,0.68618) (-5.30606,0.69943) (-4.07121,0.71758) (-3.17054,0.74234) (-2.51273,0.77615) (-2.03262,0.82237) (-1.84382,0.85153) (-1.68332,0.88554) (-1.54654,0.92539) (-1.43071,0.97186) (-1.33257,1.02624) (-1.24997,1.08974) (-1.18076,1.16397) (-1.12317,1.25091) (-1.07600,1.35238) (-1.03797,1.47112) (-1.00822,1.60984) (-0.98598,1.77210) (-0.97074,1.96158) (-0.96211,2.18316) (-0.96399,2.74515) (-0.99184,3.51333) (-1.04828,4.56112) (-1.13910,5.99595)};
\draw[->,DeepBlue](-1.02722,1.48617)--(-0.97791,1.80235);
\draw[dashed,DeepBlue](-7,0.65)--(-3.5,0.65);
\draw[densely dotted,DeepBlue](-0.57425,-1.04534)--(-1.31397,8.92726);
\draw[thick,Gold] plot coordinates{(-6.99787,1.47710) (-5.54657,1.49957) (-4.41721,1.52849) (-3.52818,1.56595) (-2.82644,1.61424) (-2.26418,1.67658) (-2.02326,1.71431) (-1.80405,1.75716) (-1.60292,1.80586) (-1.41674,1.86114) (-1.24233,1.92392) (-1.07653,1.99535) (-0.91714,2.07639) (-0.76111,2.16857) (-0.60596,2.27341) (-0.44932,2.39251) (-0.28859,2.52783) (-0.12142,2.68133) (0.24442,3.05418) (0.67235,3.53571) (1.19080,4.15795) (1.83356,4.96153) (2.64166,5.99765)};
\draw[->,Gold](-0.33358,2.48551)--(-0.09738,2.70141);
\draw[dashed,Gold](-7,1.4)--(-3.5,1.4);
\draw[densely dotted,Gold](-1.72000,0.04000)--(4.28000,8.04000);
\draw[thick,TextGray] plot coordinates{(-6.99797,2.73919) (-5.76126,2.76997) (-4.72863,2.80748) (-3.85432,2.85342) (-3.10825,2.90940) (-2.45772,2.97782) (-1.87825,3.06127) (-1.34586,3.16315) (-0.83873,3.28768) (-0.33753,3.43963) (0.17845,3.62526) (0.73003,3.85211) (1.33787,4.12866) (2.02735,4.46640) (2.82563,4.87862) (3.76605,5.38259) (4.88473,5.99766)};
\draw[->,TextGray](0.45243,3.73374)--(0.74704,3.85865);
\draw[dashed,TextGray](-7,2.6)--(-3.5,2.6);
\draw[densely dotted,TextGray](-2.17021,1.74648)--(6.47807,6.76717);
\end{scope}\fill(0,0)circle(4pt)node[below right]{$+Ze$};
\draw[<->,Gold](-5.2,0)--node[left]{$b$}(-5.2,1.4);\draw[dashed](-5.4,0)--(0,0);
\draw[->,DeepBlue,thick](-6.5,-.7)--(-4.5,-.7)node[midway,below]{$\vect v_\infty$};
\begin{scope}[shift={(2,-1.7)}]\draw[->,Gold](0,0)--(2.5,0);\draw[->,Gold](0,0)--({2.5*cos(53.1301)},{2.5*sin(53.1301)});\draw[Gold](1,0)arc(0:53.1301:1)node[right]{$\Theta$};\end{scope}
\node[below]at(-4.5,-2.4){\textcolor{DeepBlue}{$b=0.65$}};\node[below]at(0,-2.4){\textcolor{Gold}{$b=1.4$}};\node[below]at(4.5,-2.4){\textcolor{TextGray}{$b=2.6$}};
\node[below]at(0,-3.1){同一 $B_0=\kappa/(2E_k)=0.7$；所有渐近线与各自双曲线使用同一尺度};
\end{tikzpicture}}

\providecommand{\ThermalFigure}[3]{\begin{qualityfigure}
\begin{tikzpicture}[>=Latex,font=\small,x=.8cm,y=.8cm]#1\end{tikzpicture}
\figcap{#2}\label{#3}\end{qualityfigure}}
\newcommand{\ThQElectroSystems}{\ThermalFigure{
 \draw[thick](-1,0)--(1,0);\draw[dashed](-2,0)--(-1,0) (1,0)--(2,0);\fill(0,1.6)circle(2pt)node[above]{$P$};\draw[<->](0,0)--node[right]{$a$}(0,1.6);\node[below]at(0,-.4){直线，$\lambda$};
 \begin{scope}[shift={(5,0)}]\draw[thick](-1.5,0)--(.8,0)--(1.5,.6)--(-.8,.6)--cycle;\fill(0,1.9)circle(2pt)node[above]{$P$};\draw[dashed](0,.3)--node[right]{$a$}(0,1.9);\node[below]at(0,-.4){平面或长条，$\sigma$};\end{scope}
 \begin{scope}[shift={(10.5,0)}]\draw[thick](0,0)ellipse(1.3 and .45);\draw[dashed](0,0)--(0,2);\fill(0,2)circle(2pt)node[above]{$P$};\draw(0,2)--(1.3,0);\node[right]at(1.3,1){$\theta$};\node[below]at(0,-.8){圆盘，$\sigma$};\end{scope}
}{积分对象分别为线电荷、面电荷与圆盘；有限长度的端点角约定另见题中详图。}{ch19:quality:systems-q}}
\newcommand{\ThQUnequalCharges}{\ThermalFigure{
 \draw[->](-4,0)--(4,0)node[right]{$x$};\fill[DeepBlue](-1.8,0)circle(4pt)node[above]{$+Q$};\fill(1.8,0)circle(3pt)node[above]{$-q$};
 \node[below]at(-1.8,0){$-a_0$};\node[below]at(1.8,0){$a_0$};\node[above]at(0,1.1){$Q>q>0$};
}{两电荷的位置与大小关系已知，场线、零场点及远场性质待求。}{ch19:quality:unequal-q}}
\newcommand{\ThQCavity}{\ThermalFigure{
 \fill[DeepBlue!10](0,0)circle(2);\draw[thick](0,0)circle(2);\fill[white](.7,.5)circle(.8);\draw[thick](.7,.5)circle(.8);
 \fill(0,0)circle(1.5pt)node[below]{$O$};\fill(.7,.5)circle(1.5pt)node[above]{$C$};
 \node[right,align=left]at(3.2,.7){实心部分固定体电荷密度 $\rho$\\空洞内无电荷\\另比较静电平衡金属};
}{偏心球形空洞；绝缘固定体电荷与导体感应电荷是两种不同模型。}{ch19:quality:cavity-q}}
\newcommand{\ThQPotential}{\ThermalFigure{
 \draw[thick](-2,0)--(2,0);\node[below]at(0,-.2){线电荷 $\lambda$};
 \fill(0,1.5)circle(2pt)node[right]{$P$};\fill(0,3)circle(2pt)node[right]{$P_0$};\draw[<->](-.5,1.5)--node[left]{$\Delta s$}(-.5,3);
 \node[right,align=left]at(3.5,2){比较电势与电势差的积分\\有限和无限带电体分别检查收敛};
}{用两个有限位置说明电势差的选取；其余对称系统沿用各自的题设几何。}{ch19:quality:potential-q}}
\newcommand{\ThQZeroPotential}{\ThermalFigure{
 \draw[->](-1,0)--(6,0)node[right]{$x$};\fill(0,0)circle(3pt)node[below]{$-q$};\fill[DeepBlue](3,0)circle(3pt)node[below]{$+Q$};
 \fill(1.5,2)circle(2pt)node[above]{$P(x,y,z)$};\draw[dashed](0,0)--node[left]{$r_-$}(1.5,2);\draw[dashed](3,0)--node[right]{$r_+$}(1.5,2);
 \node[below]at(3,-.5){$x=a$};
}{零等势面上的一般点到两个电荷的距离不同；图未预设球心或半径。}{ch19:quality:zero-q}}
\newcommand{\ThZeroPotentialAnalysisDrawing}{\begin{tikzpicture}[>=Latex,font=\small,x=.9cm,y=.9cm]
 \draw[->](-4,0)--(4,0)node[right]{$x$};\draw[DeepBlue,thick](-1,0)circle(2);
 \fill(0,0)circle(3pt)node[below]{$-q$};\fill[DeepBlue](3,0)circle(3pt)node[below]{$+Q$};
 \fill(-1,0)circle(1.5pt)node[below]{$x_c$};\draw[<->](-3,2.5)--node[above]{$2R_0$}(1,2.5);
 \fill(0,1.73205)circle(1.5pt)node[above]{$P$};\draw[dashed](0,0)--(0,1.73205);\draw[dashed](3,0)--(0,1.73205);
 \node[right,align=left]at(4.2,1.3){图示 $Q>q$ 的一种比例\\$R_0=a\mu/|\mu^2-1|$\\$x_c=-a/(\mu^2-1)$};
\end{tikzpicture}}
\newcommand{\ThQTwoLines}{\ThermalFigure{
 \draw[->](-4,0)--(4,0)node[right]{$x$};\draw[->](0,-.5)--(0,3)node[above]{$y$};
 \draw[thick](-2,0)circle(.1);\draw[thick](2,0)circle(.1);\node[below]at(-2,-.2){$-a,-\lambda$};\node[below]at(2,-.2){$a,+\lambda$};
 \fill(.7,1.8)circle(2pt)node[above]{$P$};\draw[dashed](-2,0)--(.7,1.8);\draw[dashed](2,0)--(.7,1.8);
 \node[right]at(4.5,1.5){两导线均垂直纸面，沿 $z$ 延伸};
}{在垂直导线的截面研究等势线与场线，距离均指到对应导线的垂距。}{ch19:quality:lines-q}}
\newcommand{\ThQFiniteLine}{\ThermalFigure{
 \draw[->](-4,0)--(4,0)node[right]{$x$};\draw[line width=2pt](-2,0)--(2,0);\node[below]at(-2,0){$-a$};\node[below]at(2,0){$a$};
 \fill(.7,2)circle(2pt)node[above]{$P$};\draw[dashed](-2,0)--node[left]{$r_-$}(.7,2);\draw[dashed](2,0)--node[right]{$r_+$}(.7,2);
 \node[below]at(0,-.7){线密度 $\lambda$，绕 $x$ 轴旋转对称};
}{有限直线的两个端点及一般观察点；不能预设场线为同心圆的径向线。}{ch19:quality:finite-q}}
\newcommand{\ThQElectronLens}{\ThermalFigure{
 \draw[->](-4,1.5)--(-1,1.5)node[midway,above]{入射电子};\draw[dashed](-1,1.5)--(4,1.5);
 \draw[thick](0,0)circle(.12);\node[below]at(0,-.2){$+\lambda$，金属线垂直纸面};
 \draw[<->](.8,0)--node[right]{$b$}(.8,1.5);\node[above]at(0,2.8){先经 $U_0$ 加速，偏转为小量};
}{瞄准参数从未受扰动的入射直线量到带电细线。}{ch19:quality:lens-q}}
\newcommand{\ThQSphereEnergy}{\ThermalFigure{
 \foreach\s/\lab in {0/金属球壳,5/绝缘实心球}{\begin{scope}[shift={(\s,0)}]\draw[thick](0,0)circle(1.3);\draw[->](0,0)--(45:1.3)node[right]{$R$};\node[below]at(0,-1.7){\lab，$Q$};\end{scope}}
}{比较同一外半径、同一总电荷的表面分布与体分布。}{ch19:quality:sphere-q}}
\newcommand{\ThQAssembly}{\ThermalFigure{
 \foreach\x/\y/\i in {0/0/1,1.5/.5/2,3/-.2/3,4.5/1/4,6/-.1/5,.6/2/6,2.2/2.2/7,3.9/2.6/8,5.7/2/9,7.2/1/10}{\fill[DeepBlue](\x,\y)circle(2pt)node[above]{$Q_{\i}$};}
 \draw[dashed](0,0)--node[below]{$r_{13}$}(3,-.2);\draw[dashed](1.5,.5)--node[above]{$r_{29}$}(5.7,2);
}{十个点电荷的固定末位置；距离 $r_{ij}$ 对应无序电荷对 $i,j$。}{ch19:quality:assembly-q}}
\newcommand{\ThQBohr}{\ThermalFigure{
 \draw[thick](0,0)circle(1.8);\fill[DeepBlue](0,0)circle(4pt)node[below]{$+e$};\fill(1.8,0)circle(3pt)node[right]{$-e$};
 \draw[->,thick](1.8,0)--(1.8,1.3)node[right]{$v=c/137$};\draw[<->](0,-.5)--node[below]{$r$}(1.8,-.5);
}{玻尔模型的圆轨道与给定速度；半径和总能量尚待动力学条件确定。}{ch19:quality:bohr-q}}
\newcommand{\ThQRutherford}{\ThermalFigure{
 \fill[DeepBlue](0,0)circle(5pt)node[below]{$+Ze$（固定）};\draw[->,thick](-4,1.6)--(-1.6,1.6)node[midway,above]{$\alpha:\ +2e,\ v_\infty$};
 \draw[dashed](-4,0)--(0,0);\draw[<->](-3.8,0)--node[left]{$b$}(-3.8,1.6);\node[right]at(2,1.4){不同 $b$ 对应不同入射条件};
}{固定金核与远处入射束；轨迹和散射角在解答中由能量、角动量确定。}{ch19:quality:rutherford-q}}
